\section{On Lacking Credentials}

You know when the living thought of a writer becomes fossilized when it turns into text for use by academics, particularly those who find it necessary to boast about their credentials\footnote{\url{https://gornahoor.net/?p=369}} before speaking.

I would like to remind visitors that Evola refused to accept a degree in engineering, because he did not want approval from the system. He had no credentials in philosophy, religion or politics. His German was self-taught, in order to read 19th century German philosophy.

He had no expertise in Asian languages. He read Hindu and Buddhist texts in translations. He even commented on them, apparently something an academic would not dare do. Neither Guenon nor Evola would be considered for an academic post anywhere, nor would their writings meet the standards for inclusion in academic journals.

The question about Guenon's conversion to Islam is dealt with in \textit{Did Guenon convert?}\footnote{\url{https://gornahoor.net/?p=262}}. He did not believe that there was a valid initiatory tradition in the West; that is why he became a Sufi.

The difference in Guenon's and Evola's relationship to current political issues is dealt with here\footnote{\url{https://gornahoor.net/?p=273}}. This does not mean that Guenon was effeminate, nor that his positions do not have political consequences.

The Yoga that Evola wrote about was based on the Hindu Tantric writings, not what is commonly made available in the West. Guenon wrote about the Vedas and did not regard the Tantras as orthodox.

Evola considered ancient India to be an Aryan civilization, hence he considered their religious scriptures as part of Western heritage. Nowadays, the so-called Aryan Invasion Theory is not so popular (among ``academics''), but both Guenon and Evola accepted it as true. But the real question is: are the spiritual teachings of ancient Buddhism, the Vedas, and the Tantras alien to the Western spirit or not? Of course, this presupposes that one understands the Western spirit. That cannot be learned from books, but only through one's own personal spiritual development.

Evola never advocated a return to the ancient pagan beliefs of Europe, whether Greek, Roman, or Germanic. It was more a state of mind that he endorsed, particularly as found in the mystery cults. Furthermore, he never promoted ``polytheism'', a metaphysical absurdity, if it implies the denial of a single fundamental reality. Rather his thought ran more along the lines of Herman Wirth\footnote{\url{https://gornahoor.net/library/UltimaThule.pdf}}. His position was actually closer to atheism, as God could only be realized in the Absolute man.

I remind everyone that to be ``Traditional'' requires a direct knowledge and understanding of metaphysical principles, and this comes not from academic achievement, but rather from self-development. Although the applications of principles to concrete situations may be subjects for discussion, the principles themselves are immutable and not open to debate.


\flrightit{Posted on 2010-04-13 by Cologero}

\begin{center}* * *\end{center}

\begin{footnotesize}
\begin{sffamily}
\texttt{Dart on 2015-12-19 at 12:48 said:}

``Guenon wrote about the Vedas and did not regard the Tantras as orthodox.''

In ``Studies in Hinduism,'' (chapters 7 \& 8) Guenon refers to Tantra as an orthodox viewpoint, and even says it is particularly suited to the Kali Yuga. But he warns of falling into the error of seeing magic as an end in itself rather than as a support for realization (as does Evola iirc).

\hfill

\texttt{Cologero on 2020-04-13 at 20:01 said:}

Thanks, Dart. Guenon views the Tantras as an extension of the Vedas. BTW, Wolfgang Smith also views the Tantras as orthodox.

\end{sffamily}
\end{footnotesize}

